%% ma: L12-B14-0101
%% lop: 12
%% chuong: 5
%% bai: 14
%% dang: 12.14.37
%% loai: TLN
%% muc: TH
%% dap_an: "4,5"
%% nguon: "Ảnh giáo viên gửi 07/11/2026 (Tư Duy Mở - tuduymo.com, đề tự luyện Phương trình mặt phẳng, đề lập phương trình mặt phẳng), Câu 19"
%% kien_thuc: "Bài 14 (phương trình mặt phẳng: điểm thuộc mặt phẳng), giải hệ ba phương trình bậc nhất ba ẩn (lớp 10)"
%% trang_thai: chinh-thuc
%% ghi_chu: "Tài liệu gốc không kèm đáp án; đáp án do Claude giải và kiểm bằng sympy. Có thể dùng máy tính cầm tay giải hệ. Dạng đã duyệt 07/11/2026."
\begin{ex}
Trong không gian $Oxyz$, cho ba mặt phẳng $(\alpha)\colon x-y+2=0$, $(\beta)\colon y+2z-6=0$, $(\gamma)\colon x+y-4z-12=0$. Gọi $M$ là điểm nằm trên cả ba mặt phẳng đã cho. Hãy xác định hoành độ của điểm $M$ (làm tròn kết quả đến hàng phần mười).
\shortans{4,5}
\loigiai{Giải hệ $\begin{cases}x-y+2=0\\y+2z-6=0\\x+y-4z-12=0\end{cases}$ được $x=\dfrac92$, $y=\dfrac{13}2$, $z=-\dfrac14$. Hoành độ $4{,}5$.}
\end{ex}
